
\begin{abstract}
Solar storms are able to cause damage to vulnerable space and ground based platform components. These solar storms consists of many classes of particles such as electrons, protons, alpha particles, heavy nuclei and neutrons. Particles that cause the most damage to equipment are protons  $>$ 10 MeV  colliding with satellites and discharging above the critical charge, $Q_c$, to cause single event upsets (SEUs). An SEU is the process in which a bit flip occurs in the static random access memory unit causing anomalies in data or soft errors. To safeguard against these events we use two  current relativistic electron intensity channels ($>.25$ and $>0.67$) and proton fluxes from solar storms to forecast future proton intensity levels. This forecasting method was first introduced using a forecasting matrix mapping the Electron Increase Parameter to the electron intensity which would give a heatmap of the proton intensity. With the advent of machine learning, Neural Networks and Recurrent Neural Networks provided better forecasting F1 values than it predecessors. We choose a collection of transformer architectures to evaluate against each previous models and compare against Mean Absolute Error (MAE), lag from onset to peak fluxes, and F1 scores. We demonstrate a reduction in the MAE across all transformer architectures with similar lag values. On average we find improved F1 scores marking a well designed forecasting model. This paper will demonstrate a baseline case to make related predictions on other satellite platforms such as GOES.  
\end{abstract}

\clearpage

\section{Introduction}
Space based equipment and platforms are vulnerable to their surrounding Space Weather Environment (SWE) . The Sun dominates the SWE with solar storms comprised of winds that travel from its surface to us. The storms themselves are made from a composite of classes of particles such as protons, electrons, alpha particles, and heavier nucleons. Together, because of their origination, they make up what we call Solar Energetic Particles (SEP)s. The SEP's travel from our sun and collide with our satellites and equipment resulting in damage done to the Static Random Access Memory (SRAM) of their components \cite{noeldeke2021single}. Because of their abundance and difficulty in shielding,  high energy protons ($>10$ MeV) are the largest culprit of damaging our equipment. \cite{torres2025machine}.     The protons cause "soft" anomalies in the Static Random Access Memory (SRAM) resulting in bit flips in the telemetry data \cite{carlton2018telemetry} if the energy of the proton is larger than the critical charge, $Q_c$, of the equipment \cite{barak1999figure}. 

Relativistic electrons will typically arrive 1 AU ahead of non-relativistic  high energy protons due to the less scattering, making them a hallmark for predicting solar storms \cite{posner2007forecasting}. This relationship was demonstrated by \cite{vanhollebeke1975variation} $\leftarrow$ \hl{can't access}, who describing electrons as the first in situ sign of a particle event.  A forecasting matrix method mapping the electron increase parameter to the electron intensity creating a heat map for the average upcoming proton intensities for each electron rise-electron flux pair combination \cite{posner2007forecasting}. Machine learning algorithms were later developed and an Neural Network (NN) and Recurrent Neural Network (RNN) was able to achieve a higher forecasting F1 score then its  forecasting matrix predecessor \cite{torres2025machine}. 

In this paper, we test if a time series transformer is able to attain improved forecasting scores than the previously stated methods? Examine multiple positional encoder architectures to create these predictions.

 Section \ref{sec:relatedwork} will discuss related works replacing forecasting RNN models with transformers. Section \ref{sec:approach} will briefly introduce the RNN model we will be replicating along with the chosen positional encoding methods to compare forecasting accuracy.  We measure our models outputs with a persistence model as a baseline \cite{torres2025machine} and compare our error, lag and forecasting. Section \ref{sec:experimentalevaluation} will compare models and show that we have improved error reduction compared to \cite{torres2025machine} and \cite{posner2007forecasting} but similar lag and forecasting. Section \ref{sec:conclusion} will summarize and discuss future work, potential improvements and limitations.
